\section{Conclusion}
\label{sec:conclusion}

We systematically evaluated Zero-Space Projection (\ZSP) on three 
perception tasks: single-frame \MOT, multi-frame \MOT, and 
cooperative perception. Our conclusions are:

\begin{enumerate}[leftmargin=*]
    \item On \textbf{single-frame \MOT}, \ZSP achieves $99.96\%$ of 
          Hungarian's accuracy but is $20{,}000\times$ slower.
    \item On \textbf{multi-frame \MOT}, \ZSP is statistically 
          significantly worse than Hungarian ($-6.7\%$, $p=0.0001$, 
          win rate $0\%$).
    \item On \textbf{cooperative perception}, all methods achieve 
          $100\%$ accuracy, but distributed Hungarian is fastest and 
          requires zero communication.
\end{enumerate}

\textbf{Root cause.} Linear assignment problems possess a dual 
structure that enables the Hungarian algorithm to achieve exact 
optimality in polynomial time. \ZSP is an iterative approximation 
and cannot, by construction, outperform the exact algorithm.

\textbf{Implication.} The value of \ZSP---and of physics-native 
computing more broadly---is not in replacing problems already 
solved by exact algorithms. It is in solving problems that lack 
such algorithms: material simulation, physics-layer cryptography, 
high-precision sensing, and variational field solving.

\textbf{Final thought.} Knowing when \emph{not} to use a tool is 
as important as knowing when to use it. We hope this negative 
result serves as a useful boundary marker for the physics-native 
computing community.